\section{Resultados}

\subsection{Actividad 1}

Durante esta etapa se desarrolló la conceptualización y el diseño de la metodología de trabajo para el procesamiento de señales EEG, orientada a la clasificación del trastorno por déficit de atención e hiperactividad (TDAH) en población infantil. Este proceso incluyó una revisión técnica focalizada en el análisis de señales EEG, arquitecturas de aprendizaje profundo y esquemas de validación cruzada aplicados a datos neurofisiológicos, lo que permitió establecer un marco metodológico consistente y alineado con los objetivos del proyecto.

Como resultado de este análisis, se definieron y evaluaron diferentes enfoques de modelado. Inicialmente, se consideró una arquitectura Transformer linealizada basada en embeddings de Takens (LinearTakensFormer), debido a su capacidad para capturar dependencias temporales en señales no estacionarias. No obstante, evaluaciones preliminares evidenciaron limitaciones en su capacidad de representación. En consecuencia, se establecieron dos líneas de modelado principales que utilizan la misma base de datos EEG ADHD/Control Children: i) un modelo híbrido que combina EEGNet para la extracción de características espaciales y espectrales con un Transformer Encoder para el modelado de dependencias temporales de largo alcance, el cual sustenta el artículo de investigación presentado en el Anexo 1; y ii) un modelo alternativo basado en el uso de un kernel personalizado tipo RBF, cuya formulación mejora el desempeño frente a kernels estándar, documentado en el Anexo 2. Ambas aproximaciones responden directamente al objetivo de mejorar la capacidad de discriminación entre sujetos con TDAH y controles.

En esta fase también se definió la base de datos de trabajo, compuesta por registros EEG de 19 canales adquiridos a 128 Hz bajo el sistema internacional 10-20, correspondientes a sujetos en condición de reposo clasificados en dos grupos (TDAH y control). Asimismo, se estableció una cadena de preprocesamiento que incluye filtrado pasabanda entre 0.5 y 60 Hz, eliminación de interferencia de red en 48–52 Hz y segmentación mediante ventanas deslizantes con 50 \% de solapamiento, garantizando consistencia en la preparación de los datos para ambos modelos.

Finalmente, se diseñó el esquema de validación basado en separación por sujetos bajo el enfoque Leave-K-Subjects-Out, implementado mediante validación cruzada estratificada tipo nested group k-fold con $k=5$, asegurando evaluación generalizable y evitando fuga de información.

\subsection{Actividad 2}

En esta etapa, actualmente en desarrollo, se ha llevado a cabo una revisión sistemática del estado del arte en el análisis de señales EEG, con énfasis en técnicas de aprendizaje profundo y enfoques híbridos para clasificación en contextos clínicos. Esta revisión ha permitido validar las decisiones metodológicas adoptadas en la Actividad 1, particularmente en relación con el uso de arquitecturas que integran información espacial, espectral y temporal, así como el uso de esquemas de validación robustos basados en separación por sujetos.

Los resultados preliminares evidencian que las arquitecturas híbridas, como la combinación EEGNet y Transformer, presentan ventajas significativas en la modelación de señales EEG complejas, lo cual es consistente con los resultados hasta el momento en el desarrollo experimental asociado al Anexo 1. De manera complementaria, la exploración de kernels personalizados, como el RBF propuesto, se alinea con tendencias recientes orientadas a mejorar la capacidad de generalización en espacios de características no lineales, en concordancia con los avances documentados en el Anexo 2.

\subsection{Actividad 3}

En el marco del desarrollo e implementación de modelos de análisis de señales EEG, durante esta etapa se consolidó el modelo propuesto denominado \textit{EEG-TACT: Temporal Attention and Convolutional Tokenization for Interpretable ADHD Classification from EEG}, integrando los componentes metodológicos definidos en fases previas y validando experimentalmente su desempeño bajo un esquema de evaluación independiente por sujetos. La arquitectura desarrollada combina un bloque convolucional inspirado en EEGNet para la extracción de representaciones espaciales y espectrales, un operador Transformer Encoder para el modelado de dependencias temporales y un mecanismo de \textit{attention pooling} orientado a mejorar la interpretabilidad del modelo y la agregación de información relevante en la secuencia temporal.

Como resultado del proceso de validación y optimización, el modelo alcanzó un desempeño estable y competitivo en la tarea de clasificación TDAH/control, obteniendo una exactitud a nivel de sujeto de 87.5 \%, \textit{recall} de 96.0 \% y precisión de 82.8 \%, manteniendo además un número reducido de parámetros entrenables. De manera complementaria, se realizaron análisis de interpretabilidad, estudios de ablación y comparaciones estadísticas frente a modelos de referencia reportados en la literatura, evidenciando mejoras en robustez y consistencia del etiquetado bajo diferentes inicializaciones experimentales. Estos resultados dieron lugar a la preparación, validación y posterior envío del artículo científico asociado a una revista especializada.
